%! AP Statistics — Unit 8, Lesson 8.3 — Group Activity
\documentclass[10pt]{article}
\usepackage{apstats-article}
\usepackage{apstats-boxes}

\begin{document}

\pageheader{Unit 8, Lesson 8.3}{Group Activity}
\namepartnerperiod

\vspace{0.05in}

\begin{tcolorbox}[colback=sky!30, colframe=navy, boxrule=0.8pt, arc=2mm,
    left=3mm, right=3mm, top=2mm, bottom=2mm]
\textbf{Birth Month Distribution of MLB Players}\\[4pt]
Researchers wonder whether Major League Baseball (MLB) players are more likely to be born
in certain months. If birth month had no effect on player development, we would expect
players to be born equally often in each month ($p = 1/12 \approx 0.0833$ for each month).
A random sample of 180 MLB players yields the following birth month counts:

\renewcommand{\arraystretch}{1.4}
\begin{center}
\begin{tabular}{lc|lc}
    \textbf{Month} & \textbf{Observed} & \textbf{Month} & \textbf{Observed} \\
    \hline
    January   & 21 & July      & 13 \\
    February  & 18 & August    & 12 \\
    March     & 19 & September & 16 \\
    April     & 22 & October   & 14 \\
    May       & 17 & November  & 11 \\
    June      & 15 & December  & 22 \\
\end{tabular}
\end{center}
\end{tcolorbox}

\vspace{0.08in}

\begin{enumerate}[leftmargin=1.5em]

  \item \textbf{Hypotheses.} State $H_0$ and $H_a$ for a chi-square goodness-of-fit test.

        \vspace{0.08in}
        $H_0$: \blank{10cm}

        \vspace{0.12in}
        $H_a$: \blank{10cm}

        \vspace{0.08in}

  \item \textbf{Expected Counts.} Calculate the expected count for each month and verify
        the large-counts condition.

        \vspace{0.05in}
        $E = n \times (1/12) =$ \blank{2cm} $\times$ \blank{2cm} $=$ \blank{2cm}
        for each month.

        \vspace{0.05in}
        Large-counts condition: \blank{10cm}

        \vspace{0.08in}

  \item \textbf{Test Statistic.} Compute the $\chi^2$ statistic. Fill in each contribution
        below, then sum.

\renewcommand{\arraystretch}{2.0}
\begin{center}
\begin{tabular}{lcccc}
    \textbf{Month} & $O$ & $E$ & $(O-E)^2/E$ & \\
    \hline
    January   & 21 & \blank{1.2cm} & \blank{2cm} & \\
    February  & 18 & \blank{1.2cm} & \blank{2cm} & \\
    March     & 19 & \blank{1.2cm} & \blank{2cm} & \\
    April     & 22 & \blank{1.2cm} & \blank{2cm} & \\
    May       & 17 & \blank{1.2cm} & \blank{2cm} & \\
    June      & 15 & \blank{1.2cm} & \blank{2cm} & \\
    July      & 13 & \blank{1.2cm} & \blank{2cm} & \\
    August    & 12 & \blank{1.2cm} & \blank{2cm} & \\
    September & 16 & \blank{1.2cm} & \blank{2cm} & \\
    October   & 14 & \blank{1.2cm} & \blank{2cm} & \\
    November  & 11 & \blank{1.2cm} & \blank{2cm} & \\
    December  & 22 & \blank{1.2cm} & \blank{2cm} & \\
    \hline
    \textbf{Total} & 180 & & \blank{2.5cm} $= \chi^2$ & \\
\end{tabular}
\end{center}

  \item \textbf{Degrees of Freedom and P-Value.}

        $df =$ \blank{3cm}

        \vspace{0.05in}
        Using a chi-square table with $df = 11$, the p-value is approximately 0.30.
        \textbf{Interpret the p-value} in context.

        \writelines{3}

  \item \textbf{Conclusion.} State a complete conclusion at $\alpha = 0.05$.

        \writelines{3}

\end{enumerate}

\end{document}
